% Generated draft: requires mathematical and source review.
\documentclass[11pt]{article}
\usepackage[margin=25mm]{geometry}
\usepackage{fontspec}
\setmainfont{Latin Modern Roman}
\usepackage{amsmath,amssymb,mathtools,unicode-math}
\setmathfont{Latin Modern Math}
\usepackage{enumitem,needspace,xurl,hyperref}
\setlength{\parindent}{0pt}
\setlength{\parskip}{5pt}
\setlength{\emergencystretch}{3em}
\newcommand{\N}{\mathbb{N}}
\newcommand{\Z}{\mathbb{Z}}
\newcommand{\Q}{\mathbb{Q}}
\newcommand{\R}{\mathbb{R}}
\newcommand{\C}{\mathbb{C}}
\newcommand{\F}{\mathbb{F}}
\newcommand{\A}{\mathbb{A}}
\newcommand{\PP}{\mathbb{P}}
\newcommand{\E}{\mathbb{E}}
\newcommand{\eps}{\varepsilon}
\newcommand{\dd}{\,\mathrm{d}}
\begin{document}
\section*{UnsolvedMath ID: 10400201}
\textbf{Source problem code:} catalogue entry
\subsection{Definitions and mathematical statement}
Let \(Z_{\mathrm{LMO}}\) denote the source's Le--Murakami--Ohtsuki invariant of three-manifolds. The task is to construct its analogue with coefficients in a finite field \(\F_q\), where \(q\) is a prime power. The intended field or quantifier over fields is not stated.

For the usual closed Jacobi-diagram target, one possible coefficient framework is
\[\widehat{\mathcal A}(\varnothing;\F_q)=\prod_{d\geq0}\mathcal A_d(\varnothing;\F_q),\]
where \(\mathcal A_d\) is spanned by vertex-oriented trivalent diagrams of degree \(d\) (half the number of vertices), modulo the specified antisymmetry and IHX relations. A finite-field construction in this framework would give \(Z_{\mathrm{LMO},\F_q}(M)\) in this completion, invariant under the appropriate equivalences of admissible three-manifolds and with a specified relation to the original LMO theory. This is a conditional target description; the source does not reproduce its diagram or normalization conventions or say whether an integral, truncated or modified target is intended.

The phrase ``with coefficients in a finite field'' cannot mean reducing arbitrary rational coefficients without further work. If \(\operatorname{char}\F_q=p\), there is no unital ring map \(\Q\to\F_q\), since the invertible rational integer \(p\) would map to zero. One therefore needs a stated integral or localized coefficient model with valid reduction, a degree/characteristic restriction, or a different construction. The source does not select among these possibilities or restrict the degree. It likewise does not specify whether all closed three-manifolds, rational homology spheres or integral homology spheres are intended. An arbitrary finite-field-valued topological invariant alone would not identify the requested LMO analogue.
\subsection{Short English statement}
Construct a finite-field version of the LMO invariant, specifying its diagram target and compatibility with the original theory. Rational coefficients cannot simply be reduced modulo a prime without a valid coefficient model.
\subsection{Sources}
\begin{itemize}
\item Local UnsolvedMath registry, permanent ID \(10400201\), exact source code \texttt{catalogue entry}.
\item Ohtsuki, \emph{Problems on invariants of knots and 3-manifolds} (2002), Problem 11.5, PDF page 156: \url{https://msp.org/gtm/2002/04/gtm-2002-04-024s.pdf}. Registry access 2026-07-29; extraction pdf-font-structure.
\end{itemize}
\end{document}
